\section{Methodology}\label{sec:method}

\subsection{CPFC Protocol Overview}

The CPFC protocol defines four necessary conditions:
\begin{itemize}
    \item \textbf{P1 - Discrete Checkpointability}: Tasks must decompose into discrete checkpoints
    \item \textbf{P2 - External Rule Availability}: External rules must exist for verification
    \item \textbf{P3 - Ordinal Categorical Reducibility}: Failures must form ordered categories
    \item \textbf{P4 - Post-Hoc Availability}: Outcomes must be observable post-generation
\end{itemize}

We emphasize that these conditions are necessary but not sufficient for general validity. The protocol constrains applicability; it does not establish architecture-independent generalization.

\subsection{Verification-Entropy as Ordinal Diagnostic}

The verification-entropy $H_v$ is an ordinal index computed from checkpoint-frequency tables. It is not a cardinal uncertainty measure. The appropriate interpretation is within-protocol diagnostic, not cross-domain comparison.

The $R^2$ values reported in our analysis are diagnostic indicators of ordinal pattern stability. They do not represent explanatory power in the statistical sense, nor do they prove causal relationships. We report $R^2$ with the explicit caveat that it indicates only the consistency of ordinal relationships within the protocol.
